%! Carrying Out a Test for the Difference of Two Population Proportions — Unit 6, Lesson 6.11 — Group Activity
\documentclass[10pt]{article}
\usepackage{apstats-article}
\usepackage{apstats-boxes}
\newcommand{\TallMath}[1]{$\displaystyle #1\rule[-1.4em]{0pt}{3.2em}$}

\begin{document}

\pageheader{Unit 6, Lesson 6.11}{Group Activity}
\namepartnerperiod

\vspace{0.15in}

\textbf{Directions:} Complete the full four-step procedure for each scenario. Use $\alpha = 0.05$.

\begin{scenariobox}[Scenario 1 --- Volunteering Rates]{sky}
A nonprofit wants to know whether the proportion of adults who volunteered in the past year
differs between two cities. City M: 78 of 260 volunteered. City N: 84 of 210 volunteered.

\begin{enumerate}[label=\arabic*.]
  \item \textit{State:} Define parameters, state $H_0$ and $H_a$. \writelines{2}

  \item \textit{Plan:} Compute $\hat p_M$, $\hat p_N$, $\hat p_c$. Verify conditions.
        \writelines{4}

  \item \textit{Do:} Compute $z$ and the $p$-value.
        \writelines{3}

  \item \textit{Conclude:} Write a complete conclusion. \writelines{2}
\end{enumerate}
\end{scenariobox}

\begin{scenariobox}[Scenario 2 --- Flu Shot Acceptance]{sky}
A public health study tests whether adults aged 18--34 are less likely to get a flu shot
than adults aged 55+. Sample of 18--34 year olds: 63 of 210 got a shot.
Sample of 55+ year olds: 112 of 280 got a shot.

\begin{enumerate}[label=\arabic*.]
  \item \textit{State:} Define parameters and state the one-sided hypotheses. \writelines{2}

  \item \textit{Plan:} Compute $\hat p_c$. Verify conditions.
        \writelines{3}

  \item \textit{Do:} Compute $z$ and the $p$-value.
        \writelines{3}

  \item \textit{Conclude.} \writelines{2}
\end{enumerate}
\end{scenariobox}

\end{document}
